\section{Možná rozšíření}
Ačkoliv implementovaná aplikace splňuje veškeré požadavky definované v zadání diplomové práce, během procesu vývoje byly identifikovány oblasti, které by mohly být předmětem dalšího rozvoje. Pro zajištění efektivnější práce s rozsáhlejším množstvím dat se jako vhodné jeví rozšíření implementace o mechanismy řazení a filtrování v rámci tabulkových přehledů. Tímto způsobem by uživatelům bylo umožněno dynamicky organizovat zobrazené výstupy na základě specifických podmínek.

V rámci zvýšení uživatelské přívětivosti hlavního dashboardu by mohly být implementovány přímé odkazy do zdrojových systémů GitHub a GitLab. Jednotlivé záznamy v tabulkách by tak mohly obsahovat hypertextové odkazy směřující na detail konkrétního pull requestu v externím systému, čímž by byla zajištěna hlubší provázanost aplikace s verzovacími systémy. Kromě zmíněné integrace by dashboard mohl být vylepšen o perzistenci stavu filtrů (pro uživatele a časové období), které jsou v aktuálním řešení resetovány po~opuštění stránky. Implementací mechanismu pro ukládání jejich stavu by byla eliminována nutnost opakované konfigurace pohledů po každém přihlášení či přechodu mezi stránkami systému.

Vzhledem k tomu, že aplikace agreguje a normalizuje data z různých systémů pro správu verzí softwaru, disponuje značným množstvím informací s vysokou vypovídající hodnotou. Tato data mají vysoký potenciál pro další analytické využití, a proto je jako další možné rozšíření navrhováno napojení na~externí nástroje business intelligence. Data aplikace by následně mohla sloužit pro tvorbu komplexních reportů či hlubší analýzu trendů, kterou implementovaný systém v současné době nepokrývá.
